\chapter{Orchestration and Reliability}
\label{chapter:Orchestration}
% EVOLVE-BLOCK-START

A single agent executing a single task is the simplest case of agentic work.\index{Orchestration}
Some research tasks are more demanding: they require long, multi-step computations; they benefit from having one agent check another's work; they must run reliably for hours without supervision; and they must report failures honestly when they do fail.
\emph{Orchestration} is the practice of composing multiple agents and operations into a trustworthy workflow.

The orchestrated pipeline built in this course previews the evolutionary framework.
The loop implemented here — \textbf{propose → evaluate → select → repeat} — is the same loop that becomes an evolutionary algorithm when a mutation operator is added at the pivot.
Section~\ref{section:MultiStepWorkflows} composes single operations into pipelines; Section~\ref{section:MultiAgentArchitectures} adds a second agent to check the first; Section~\ref{section:MemoryState} handles what must be remembered across steps; Section~\ref{section:Reliability} makes the pipeline survive long unattended runs; and Section~\ref{section:GenEvalSel} distills all of it into the generate--evaluate--select loop that the next chapter turns into evolution.

\section{Multi-Step Workflows}
\label{section:MultiStepWorkflows}

A \emph{workflow}\index{Workflow} is a directed graph of operations: some operations depend on the outputs of others; some can run in parallel.
The orchestrator's job is to schedule these operations correctly and collect their results.
The graph is a picture of dependency, not of time: an edge means the downstream operation cannot start until the upstream one has produced what it needs, while two operations with no path between them may run in either order, or at once.
A study that loads a dataset, fits three candidate models to it, and then compares their errors is already a graph of this shape: the three fits each depend on the load, none depends on another, and the comparison depends on all three.

\subsection{Pipeline Composition}

In the simplest case, a workflow is a \emph{pipeline}\index{Pipeline}: a linear sequence of operations where the output of each step is the input to the next.

\begin{example}
A three-step research pipeline:
\begin{lstlisting}[language=python]
def run_pipeline(dataset_path, k=5):
    # Step 1: retrieve domain knowledge
    context = retriever.query(dataset_path, k=k)

    # Step 2: agent proposes and scores candidates
    candidates = proposer.run(dataset_path, context=context)

    # Step 3: select the best candidate
    best = min(candidates, key=lambda c: c.score)
    return best
\end{lstlisting}
Each step can be a model call, a tool invocation, or a subprocess; the pipeline connects them.
Read the data flow rather than the syntax: \code{context} is produced by step~1 and consumed by step~2, and \code{candidates} is produced by step~2 and consumed by step~3, so these three lines admit only one order.
Note also that \code{min} is the correct selector only because the score here is an error, for which lower is better (the course-wide convention; Section~\ref{subsection:MultiObjectiveScoring}); under a fitness convention where higher is better, the same line would need \code{max}.
\end{example}

\subsection{Parallel Execution}

Some operations are independent: multiple agents can propose candidates simultaneously, or a critic can score candidates in parallel with new proposals being generated.
\emph{Parallel execution}\index{Parallel!execution} reduces wall-clock time by running independent operations concurrently.
The saving is largest exactly where research pipelines spend their time: if one proposer call waits ninety seconds for the provider to answer, four such calls in sequence cost six minutes and four at once cost about ninety seconds, because the waiting overlaps.

\begin{example}
Spawning multiple proposer agents in parallel (conceptual):
\begin{lstlisting}[language=python]
import concurrent.futures

with concurrent.futures.ThreadPoolExecutor(max_workers=4) as pool:
    futures = [pool.submit(proposer.run, dataset_path, seed=i)
               for i in range(4)]
    all_candidates = [f.result() for f in futures]
\end{lstlisting}
The four proposers run concurrently; the orchestrator collects and merges their outputs.
\end{example}

Four pieces of that listing behave unlike ordinary Python lines, and each is worth naming.
A \emph{thread pool}\index{Thread pool} is a small crew of workers that the program hands jobs to, and \code{max\_workers} caps how many run at once: twenty submitted jobs would still occupy at most four workers at a time, the rest queueing.
\code{pool.submit} does not run \code{proposer.run}; it queues that call and returns immediately with a \emph{future}\index{Future}, a placeholder object standing for a result that does not exist yet, which is why all four submissions complete in a fraction of a second even though no proposer has answered yet.
\code{f.result()} is where the program actually waits: it blocks until that particular proposer has finished, then hands back its return value, so by the time the comprehension on the last line ends, all four have returned.
It also re-raises whatever exception the worker raised, which matters for Section~\ref{section:Reliability}: a proposer that failed inside the pool surfaces at \code{result()} rather than vanishing silently.
The \code{with} block shuts the pool down on exit, including when an exception escapes it.

\begin{remark}
A thread pool is the right instrument here because this work is dominated by \emph{waiting}: each proposer sits idle on a network response for nearly its whole lifetime, and an idle thread costs almost nothing.
It is the wrong instrument for work dominated by computation written in pure Python, where the interpreter's global lock lets only one thread execute Python bytecode at a time and four threads finish no faster than one.
Numerical libraries such as NumPy release that lock inside their compiled kernels, so array-heavy evaluation does speed up across threads; a hand-written Python loop scoring candidates does not, and wants \code{ProcessPoolExecutor}, which uses separate operating-system processes and so escapes the shared lock, at the price of copying data between them.
\end{remark}

The principal hazard of parallel execution is \emph{non-determinism}\index{Non-determinism}: the order in which parallel operations complete is not predictable across runs.
If the final selection depends on ordering (e.g., in case of a tie), the result may differ between runs.
Concretely: two proposers independently arrive at the same expression with identical error, and the selector returns whichever of the two it met first.
When that encounter order is set by which worker happened to finish first, the same seed reports one expression on Monday and the other on Tuesday; neither result is wrong, which is precisely what makes the discrepancy hard to explain to a reviewer.
Always set seeds and record exactly which candidates were selected.

\section{Multi-Agent Architectures}
\label{section:MultiAgentArchitectures}

\subsection{The Proposer--Critic Pattern}

The most common multi-agent pattern in the course is the \emph{proposer--critic} pattern (Figure~\ref{figure:ProposerCritic}):\index{Proposer--critic pattern}
\begin{itemize}
\item A \emph{proposer} agent generates candidate artifacts: model expressions, hypotheses, experimental designs.
Its goal is breadth: generating diverse, plausible candidates quickly.
\item A \emph{critic} (or \emph{verifier}) agent examines each candidate for correctness, plausibility, or constraint satisfaction before the candidate is passed to the evaluator.
Its goal is precision: filtering out candidates that are obviously wrong.
\end{itemize}

\begin{figure}[htb!]
\begin{center}
%% proposer-critic.tex: the proposer-critic-evaluator architecture (Chapter 9).
%% House conceptual style.
%% _concept-style.tex: shared TikZ styles for the course's CONCEPTUAL block diagrams.
%% The leading underscore marks this as the one file in figures/ that is not a
%% figure; every other figures/*.tex holds exactly one figure body.
%%
%% \input this at the top of every conceptual figure (before \begin{tikzpicture})
%% so each such diagram speaks one visual language: rounded "block" nodes, stealth
%% arrows, and natural `<hue>!15' fills. Redefining these styles on each \input is
%% idempotent and harmless. The reference for the house parameters is
%% resources/STYLE-figures.md. Figure~\ref{figure:AgentLoop} is the exemplar.
%%
%% Scope: use for conceptual/relational diagrams (boxes-and-arrows). Do NOT use for
%% product mock-ups (the rig figures have their own shared language) or for
%% data/plot figures.
\tikzset{
  % the standard block node: geometry only — apply a `fill=<hue>!15' per node.
  conceptbox/.style={draw, rounded corners=4pt, minimum width=2.4cm,
                     minimum height=0.85cm, align=center, font=\small},
  % a flow arrow between blocks.
  conceptflow/.style={->, >=stealth', thick},
  % an edge/annotation label.
  conceptlabel/.style={font=\scriptsize, align=center},
}%

\begin{tikzpicture}[
    node distance=1.5cm and 1.8cm,
    box/.style={conceptbox},
    arr/.style={conceptflow}
]
\node[box, fill=blue!15]                 (prop) at (0,0) {Proposer};
\node[box, fill=orange!15, right=of prop] (crit) {Critic};
\node[box, fill=green!15,  right=of crit] (eval) {Evaluator};
\node[box, fill=gray!15,   below=of eval] (sel)  {Selector};

\draw[arr] (prop) -- node[midway, above, conceptlabel]{candidates} (crit);
\draw[arr] (crit) -- node[midway, above, conceptlabel]{survivors}  (eval);
\draw[arr] (eval) -- node[midway, right, conceptlabel]{scores}     (sel);
\draw[arr] (sel.west) -- node[midway, above, conceptlabel]{iterate / report}
    (prop.south |- sel.west) -- (prop.south);
\end{tikzpicture}%

\caption{The proposer--critic--evaluator architecture.
The proposer generates candidates, the critic filters out obvious errors, the evaluator scores the survivors, and the selector feeds the best back to the proposer for the next round.}
\label{figure:ProposerCritic}
\end{center}
\end{figure}

The critic's role is to enforce constraints that the evaluation harness does not check.
In the symbolic regression context, a critic might check:
\begin{itemize}
\item dimensional consistency of a proposed expression,
\item whether the expression is syntactically valid Python,
\item whether the expression is numerically stable on the training domain, and
\item whether the expression has already been tried in a previous round (deduplication).
\end{itemize}
Filtering before evaluation saves evaluation cost; this is especially important when evaluation requires model calls or expensive computation.
The arithmetic is worth doing once: if the critic rejects half of a hundred candidates at negligible cost and each surviving evaluation is a thirty-second simulation, the round finishes in twenty-five minutes rather than fifty.

\begin{example}
The four checks above are symbolic regression's instance of a pattern that re-skins wherever the proposer emits something executable.
An epidemiologist's critic rejects a proposed compartmental model whose compartments do not sum to the population; an economist's rejects a demand specification that slopes upward across the observed price range; an engineer's rejects a controller whose gain the actuator could not physically deliver.
Each check encodes domain knowledge the evaluation harness does not have, and the reason to house it in a separate agent rather than in the proposer's prompt is that it can then be stated, tested, and revised without disturbing how candidates are generated.
\end{example}

This proposer--critic loop is the pattern Anthropic's engineering writing (2024) names the \emph{evaluator--optimizer} workflow, and it presupposes what Chapter~\ref{chapter:Evaluation} supplies: explicit criteria against which iteration demonstrably improves the artifact, which is why the critic is a filter and not a rubber stamp.

\subsection{Verifiers}

A \emph{verifier}\index{Verifier} is an agent whose explicit task is to try to \emph{disprove} a claim or detect an error.
Where a critic checks soft constraints, a verifier checks hard correctness properties.
The distinction is best read as a difference in what a failed check licenses you to conclude.
A critic's verdict is a judgement about plausibility: a candidate it rejects was probably not worth evaluating, but it might have been, and a critic tuned too aggressively will throw away good work.
A verifier's verdict is a fact established by running something, and it is decisive: if the check below reports a negative minimum, the claim it was auditing is false, however confidently the proposer asserted it.

\begin{example}
A verifier for a data analysis step:
\begin{lstlisting}[language=markdown]
You are a verification agent. The primary agent has made the following
claim about the dataset:

  "The features are all non-negative, so log-transforms are safe."

Your job is to check this claim against the actual data. Run the
following check and report whether the claim is true or false:

  import numpy as np
  X, y = load_dataset('data/growth.npz')
  print(np.min(X), np.any(X <= 0))
\end{lstlisting}
The verifier can catch mistakes the proposer made silently, e.g., the possibility that some feature values are zero or negative.
Notice the shape of the assignment: the verifier is handed one claim and told to test it, not asked whether anything is wrong with the analysis.
A narrow, falsifiable question is one an agent can answer reliably, because there is a check that settles it; an open invitation to find problems returns plausible-sounding commentary whose accuracy nobody can assess.
\end{example}

\subsection{Dynamic Sub-Agent Spawning}

The proposer--critic topology above is \emph{static}: the agents and their wiring are fixed before the run begins.
A more flexible pattern lets an \emph{orchestrator}\index{Orchestrator agent} agent spawn specialist sub-agents on demand: calling a tool that launches a fresh agent instance with its own context window, tool set, and task, then collecting the result.
This is how a single agent decomposes a large job at run time: a research orchestrator might spawn one sub-agent per candidate dataset, or one per section of a report, each working in an isolated context so that the orchestrator's own window is not flooded with intermediate detail (Figure~\ref{figure:OrchestratorWorker}).
The reliability discipline of this chapter applies unchanged; each spawned agent can fail, so its result must be validated and its failures surfaced (Section~\ref{section:Reliability}).
Current agent SDKs (for example, Claude's agent SDK) expose this spawning as a built-in tool; it is the dynamic generalization of the fixed proposer--critic wiring.
When parallel agents must \emph{write} (not just analyze) they also need isolated working directories so their changes cannot collide; the git-worktree pattern that provides this is described in Appendix~\ref{chapter:RigReference}.

\begin{figure}[htb!]
\begin{center}
%% orchestrator-worker.tex: the orchestrator-worker pattern (Chapter 9).
%% House conceptual style; each worker is a `sub' node annotated with its own context.
%% _concept-style.tex: shared TikZ styles for the course's CONCEPTUAL block diagrams.
%% The leading underscore marks this as the one file in figures/ that is not a
%% figure; every other figures/*.tex holds exactly one figure body.
%%
%% \input this at the top of every conceptual figure (before \begin{tikzpicture})
%% so each such diagram speaks one visual language: rounded "block" nodes, stealth
%% arrows, and natural `<hue>!15' fills. Redefining these styles on each \input is
%% idempotent and harmless. The reference for the house parameters is
%% resources/STYLE-figures.md. Figure~\ref{figure:AgentLoop} is the exemplar.
%%
%% Scope: use for conceptual/relational diagrams (boxes-and-arrows). Do NOT use for
%% product mock-ups (the rig figures have their own shared language) or for
%% data/plot figures.
\tikzset{
  % the standard block node: geometry only — apply a `fill=<hue>!15' per node.
  conceptbox/.style={draw, rounded corners=4pt, minimum width=2.4cm,
                     minimum height=0.85cm, align=center, font=\small},
  % a flow arrow between blocks.
  conceptflow/.style={->, >=stealth', thick},
  % an edge/annotation label.
  conceptlabel/.style={font=\scriptsize, align=center},
}%

\begin{tikzpicture}[
    box/.style={conceptbox, minimum width=2.3cm, minimum height=1.0cm},
    sub/.style={conceptbox, minimum width=2.7cm, minimum height=0.9cm, font=\scriptsize},
    arr/.style={conceptflow}
]
\node[box, fill=orange!20] (orch) at (0,0) {Orchestrator};
\node[sub, fill=blue!12] (s1) at (4.7,1.7)  {sub-agent\\\emph{own context}};
\node[sub, fill=blue!12] (s2) at (4.7,0)    {sub-agent\\\emph{own context}};
\node[sub, fill=blue!12] (s3) at (4.7,-1.7) {sub-agent\\\emph{own context}};
\node[box, fill=green!20] (syn) at (9.4,0) {Synthesis};
\draw[arr] (orch) -- (s1);
\draw[arr] (orch) -- (s2) node[midway, above, font=\scriptsize]{task};
\draw[arr] (orch) -- (s3);
\draw[arr] (s1) -- (syn);
\draw[arr] (s2) -- (syn) node[midway, above, font=\scriptsize]{result};
\draw[arr] (s3) -- (syn);
\end{tikzpicture}%

\caption{The orchestrator--worker pattern (Anthropic, 2025).
A lead agent decomposes a task and spawns worker sub-agents that run in parallel, each with its own context window, then synthesizes their results.
Because each worker's context is isolated, the orchestrator's window is not flooded with intermediate detail and independent sub-questions are explored at once, at the cost of roughly an order of magnitude more tokens than a single agent.
The topology pays off for breadth-first tasks whose parts are genuinely independent, not for tightly coupled work.}
\label{figure:OrchestratorWorker}
\end{center}
\end{figure}

\section{Memory and State}
\label{section:MemoryState}

A single-session agent remembers its history within the context window.
An orchestrated workflow that runs many rounds, or that spawns parallel agents, needs explicit \emph{external memory}\index{External memory}: state stored outside any individual agent's context.

\subsection{What Needs to Be Remembered}

\begin{itemize}
\item \textbf{Tried expressions.} The set of model forms already evaluated, so that the proposer does not waste budget on duplicates.
\item \textbf{Best results.} The highest-scoring candidates, maintained across rounds, so that the selector can always return the best discovered so far.
\item \textbf{Failure modes.} Candidates that crashed the evaluator (zero-division, overflow), along with the error, so the proposer can avoid similar forms.
\item \textbf{Run configuration.} The random seed, model name, dataset version, and harness settings, so the run can be reproduced exactly.
\end{itemize}

\subsection{Memory in Practice}

External memory is most simply a file: a JSON or CSV record updated after each round.

\begin{lstlisting}[language=python]
import json, pathlib

class RunMemory:
    def __init__(self, path):
        self.path = pathlib.Path(path)
        self.state = json.loads(self.path.read_text()) \
            if self.path.exists() else {"tried": [], "best": None}

    def record(self, expression, score):
        self.state["tried"].append({"expression": expression,
                                    "score": score})
        if (self.state["best"] is None
                or score < self.state["best"]["score"]):
            self.state["best"] = {"expression": expression,
                                  "score": score}
        self.path.write_text(json.dumps(self.state, indent=2))

    def already_tried(self, expression):
        return any(t["expression"] == expression
                   for t in self.state["tried"])
\end{lstlisting}
Three details in that class carry the weight.
The constructor loads the file when it exists and otherwise starts an empty record, so one code path both begins a fresh run and resumes an interrupted one; nothing special has to be done to restart.
\code{record} rewrites the whole file on every call, which is wasteful and deliberate: a run killed by a cluster time limit or a dropped connection still leaves a valid, complete file on disk describing everything up to the last candidate.
\code{already\_tried} compares expression strings exactly, so \code{x0*x1} and \code{x1*x0} count as different candidates even though they are the same model, which is a real weakness and the reason the next paragraph reaches for a better index.

The memory file is the primary artifact of an orchestrated run: it records what was tried and what the best result was, enabling the run to be resumed if it is interrupted and reported honestly if it fails.

The single JSON file is the reproducible \emph{floor}, not the ceiling.
Production agents increasingly back external memory with a vector store, so that ``already tried'' can match \emph{semantically equivalent} candidates rather than only string-identical ones, and with hierarchical or episodic memory for runs that span many sessions (as of 2026).
The file-based version here is the minimal form that keeps a run resumable and honestly reportable; the richer stores are the same idea with a better index.
Two 2025 studies formalize \emph{what} an agent should keep across runs: \emph{ReasoningBank} (Ouyang et al., 2025) distills reusable reasoning strategies from past attempts, and \emph{Agentic Context Engineering} (Zhang et al., 2025) treats the agent's evolving context itself as the artifact to optimize.

\section{Reliability Engineering}
\label{section:Reliability}

An orchestrated pipeline that runs unattended for hours must be designed to fail gracefully.
Every model call can return an error; every tool can produce an unexpected result; any network connection can drop.
\emph{Reliability engineering}\index{Reliability} is the practice of anticipating these failures and building the system to handle them without losing state or producing incorrect results.

\subsection{Retries with Backoff}

The operations worth protecting here are the ones that leave the process.
A proposer round is an \emph{API call}\index{API!call} to a model provider over the network, and a critic that checks a candidate against a database or a web service issues an \emph{MCP query}\index{MCP} to a server running behind the protocol of Chapter~\ref{chapter:ContextEngineering}.
Both cross a boundary the pipeline does not control, so both can fail for reasons that have nothing to do with whether the request was well formed: the provider is momentarily rate limiting the account, a server is overloaded, a connection drops mid-response.
A \emph{retry}\index{Retry} re-attempts such a failed operation after a waiting period.
\emph{Exponential backoff}\index{Exponential backoff} increases the waiting period after each successive failure, reducing the load on the server that produced the error.

The waiting period also carries a random \emph{jitter}\index{Jitter} term, and the reason is the parallelism of Section~\ref{section:MultiStepWorkflows}.
Four proposers launched together will hit the same rate limit at the same moment; without jitter they would sleep for identical intervals and retry in lockstep, reproducing the burst that caused the error.

Retrying is only correct for \emph{transient} faults: those a second attempt could plausibly survive.
A \emph{deterministic} failure returns the same error however many times it is attempted, and the catalogue is familiar: a malformed request, an expired credential, a query naming a resource that does not exist, or a tool whose backing software is missing from the machine or installed broken.
Wrapping those in a retry loop buys nothing and costs three things: it turns a fast, informative error into a slow one, it spends tokens or quota on each doomed attempt, and it buries the actual diagnosis under a stack of timeouts.
The predicate is therefore which exceptions the loop catches, not how long it sleeps.

\begin{lstlisting}[language=python]
import time, random
from anthropic import (RateLimitError,       # 429: account is being throttled
                       APIConnectionError,   # network dropped or timed out
                       InternalServerError)  # 5xx, including 529 "overloaded"

# Transient only. A malformed request (400) or a bad key (401) is NOT here:
# those raise straight through to the caller on the first attempt.
TRANSIENT = (RateLimitError, APIConnectionError, InternalServerError)

def with_retry(fn, max_attempts=5, base_delay=1.0):
    for attempt in range(max_attempts):
        try:
            return fn()
        except TRANSIENT:
            if attempt == max_attempts - 1:
                raise                        # budget spent: surface it
            delay = base_delay * (2 ** attempt) + random.uniform(0, 1)
            time.sleep(delay)
    raise ValueError("with_retry needs max_attempts >= 1")
\end{lstlisting}

Read the listing twice, once for its control flow and once for its predicate.
\code{try} runs the call; if the call raises one of the three exception classes gathered into \code{TRANSIENT}, control jumps to the \code{except} block, the loop sleeps and comes around again, and any exception \emph{not} in that tuple propagates straight out to the caller because no handler here claims it.
That silence about the other exceptions is the design, not an omission.
The sleep doubles as \code{attempt} advances, giving delays of about one second, then two, four, and eight, and the \code{random.uniform} term is the jitter of the preceding paragraph.
The bare \code{raise} on the final attempt re-throws the exception that was caught rather than inventing a new one, which keeps the provider's own diagnosis intact for whoever later reads the traceback.

The same split governs the MCP and tool layer, though the exception types come from that client rather than from a model SDK.
An MCP server that is briefly unreachable is transient; an MCP server that is not running, or a tool that shells out to a binary the machine does not have, is deterministic, and the pipeline should say so on the first attempt rather than on the fifth.
Note also that provider SDKs already retry transient errors internally a few times, so a hand-written loop like this one is a coarser outer guard, not the first line of defense.
When the retry budget is exhausted, the loop re-raises rather than returning a fallback value, which is Principle~\ref{principle:HonestFailure} applied at the level of a single call.

\subsection{Logging}

\emph{Logging}\index{Logging} records what the pipeline did, when, and with what results.
A run without logs is a run that cannot be understood after the fact.

Every orchestrated run should log:
\begin{itemize}
\item the start time and run configuration,
\item every model call (model name, token counts, latency),
\item every tool call and its result,
\item every error and its resolution, and
\item the final result and elapsed time.
\end{itemize}

Write the log in a machine-readable form, one JSON object per line, because the questions asked of it afterwards are aggregate questions: what fraction of calls were retried, where the wall-clock time actually went, whether the best score improved monotonically across rounds.
A line carrying a timestamp, an event name, the round index, token counts, a latency, and a status field answers all three once a few thousand of its siblings are read back into a dataframe; the same information written out as English sentences answers none of them.
This is the same instinct that makes a researcher log an experiment in a table rather than a lab notebook paragraph, applied to a run whose ``experiments'' number in the thousands.

\subsection{Honest Failure Reporting}

An orchestrated pipeline fails more often than a single-step script.
The response to failure is to \emph{surface it}, not to hide it.
A pipeline that silently returns a fallback when it fails is worse than one that raises an exception: the former produces incorrect results that look correct; the latter produces no result but an informative error.

\begin{principle}
\label{principle:HonestFailure}
Report failures explicitly.
If the pipeline does not find a candidate better than the linear baseline, say so.
If three of ten proposer calls failed due to rate limits, report the fraction.
If the critic rejected every proposal in a round, report why.
A failure report is not a sign of a bad system; a failure report that is absent is a sign of an untrustworthy one.
\end{principle}

\section{The Generate--Evaluate--Select Loop}
\label{section:GenEvalSel}

The orchestration assembled in this chapter follows a pattern that recurs throughout the rest of the course.

\begin{definition}
\label{definition:GenEvalSel}
The \emph{generate–evaluate–select} (GES) loop is the procedure:
\begin{enumerate}
\item \textbf{Generate.} Produce a set of candidate artifacts $C = \{a_1, \ldots, a_m\}$.
\item \textbf{Evaluate.} Score each candidate on the reusable validation (fitness) split: $s_i = \Score(a_i, \mathcal{D}_{\mathrm{val}})$, never the sealed test split (Chapter~\ref{chapter:Evaluation}, Principle~\ref{principle:TestSplitInviolable}).
\item \textbf{Select.} Retain the top-$k$ candidates: $C^* = \mathrm{top\text{-}}k(C, \{s_i\})$.
\item Repeat with $C^*$ informing the next round of generation.
\end{enumerate}
\end{definition}

Concretely: in round~1 the proposer generates five candidate expressions; the harness scores them $0.12$, $0.08$, $0.23$, $0.07$, $0.15$; with $k=2$ (and lower scores better), the selector retains the candidates scoring $0.07$ and $0.08$ as the elites that seed round~2.
Every run of the loop is this small story repeated.
That two-item short list is the whole inheritance from one round to the next, which is what makes $k$ a consequential dial\index{Elites}: a small $k$ concentrates the next round's effort on the current leader, and a large $k$ keeps more of the field alive at the cost of pursuing weaker leads.

The \emph{val} subscript in step~2 is load-bearing rather than decorative.
The loop consults its scorer once per candidate per round, which over a full run is thousands of queries, and a split queried that often is a split the search has effectively been fit to.
The sealed test split is spent the moment the loop is allowed to see it, which is why it is scored exactly once, at the end, on the single winner.

This loop is an evolutionary algorithm in all but name.
The one ingredient missing is a \emph{mutation operator}: a procedure that derives new candidates from the survivors rather than proposing from scratch.
Adding mutation transforms the GES loop into the generate--evaluate--select--\textbf{mutate} evolutionary loop of Chapter~\ref{chapter:EvolutionaryFrameworks}.

The researcher who understands the GES loop and its failure modes (why candidates are duplicated, why the evaluator can be fooled, why the proposer collapses to a single good candidate) is ready for the evolutionary framework.
The pivot is not a new topic; it is an addition to a loop already understood.

\section*{Further Reading}

\begin{small}
\begin{enumerate}
\item Hong, S., et al., ``MetaGPT: Meta Programming for a Multi-Agent Collaborative Framework,'' \emph{ICLR}, 2024 (arXiv:2308.00352) — a framework for structured multi-agent communication.
\item Wu, Q., et al., ``AutoGen: Enabling Next-Gen LLM Applications via Multi-Agent Conversation,'' arXiv:2308.08155, 2023 — a widely used orchestration library for multi-agent conversations.
\item Chase, H., \emph{LangChain Documentation} — a survey of pipeline and memory patterns for LLM applications; available at \href{https://python.langchain.com}{python.langchain.com}.
\item \emph{A note on the ecosystem:} the multi-agent framework landscape has fragmented rapidly, so treat the AutoGen and LangChain references above as historical anchors for the conversation and pipeline patterns rather than as current installation targets. As a reasonable present-day starting point aligned with the course's default stack, Claude's agent SDK; but the field moves quickly, so check current adoption before committing to any framework (LangGraph, Swarm, and others are also in active use).
\item Workshop materials, \code{resources/claude-cli-workshop/} Lesson~04 — sub-agents, agent teams, and headless operation with Claude Code.
\item Anthropic, ``Building Effective Agents,'' \emph{Anthropic Engineering}, 2024 — the composable workflow-pattern vocabulary (prompt chaining, routing, parallelization, orchestrator--workers, evaluator--optimizer) that names the patterns of this chapter.
\item Anthropic, ``How We Built Our Multi-Agent Research System,'' \emph{Anthropic Engineering}, 2025 — a shipped orchestrator--worker system: parallel sub-agents in isolated contexts and the token-cost tradeoff of multi-agent designs.
\item Anthropic, ``Effective Harnesses for Long-Running Agents,'' \emph{Anthropic Engineering}, 2025 — session-spanning harness design and external-memory checkpointing (progress log, task list, git history) for agents that outlast a single context window.
\item Ouyang, S., et al., ``ReasoningBank: Scaling Agent Self-Evolving with Reasoning Memory,'' 2025 (arXiv:2509.25140) — distilling reusable reasoning strategies an agent carries across runs, one step past the string-matched ``already tried'' record.
\item Zhang, Q., et al., ``Agentic Context Engineering: Evolving Contexts for Self-Improving Language Models,'' 2025 (arXiv:2510.04618) — treating the agent's evolving context as the object to optimize across runs.
\end{enumerate}
\end{small}
% EVOLVE-BLOCK-END
